% ============================================================
% f2014_shanghai_science_q16_prism.tex
% 2014 上海卷（理科）第 16 题：四个棱长为 1 的正方体排成的正四棱柱
% 来源 PDF：普通高考/2014/2014上海理.pdf
% 原页：第 1 页（A3），印面 915，第 16 题题干右侧
% 底面 2×2、高 1；AB 为一条侧棱（A 在下、B 在上），
% 上底面除 B 外的八个格点为 P1—P8：
%   P1、P2 在左边（自 B 起），P3、P6 在前边，P4 为面心，
%   P5、P8 在后边，P7 在右边。
% 可见棱与上底面网格线为实线，被遮挡的棱与内部竖棱为虚线。
% 平行投影：前边方向取 1 单位 =（1.025,-0.15），竖直方向取 1 单位 =（0,-0.95）。
% ============================================================
\begin{tikzpicture}[>=Stealth, line width=0.65pt, font=\small, scale=1.1,
  line cap=round, line join=round]
  \coordinate (B) at (0,2.30);
  % 上底面九个格点
  \coordinate (P1) at ($(B)+(0.55,0.55)$);
  \coordinate (P2) at ($(B)+(1.10,1.10)$);
  \coordinate (P3) at ($(B)+(2.05,-0.30)$);
  \coordinate (P4) at ($(B)+(2.60,0.25)$);
  \coordinate (P5) at ($(B)+(3.15,0.80)$);
  \coordinate (P6) at ($(B)+(4.10,-0.60)$);
  \coordinate (P7) at ($(B)+(4.65,-0.05)$);
  \coordinate (P8) at ($(B)+(5.20,0.50)$);
  % 下底面九个格点（竖直下移 0.95）
  \coordinate (A)   at ($(B)+(0,-0.95)$);
  \coordinate (P1b) at ($(P1)+(0,-0.95)$);
  \coordinate (P2b) at ($(P2)+(0,-0.95)$);
  \coordinate (P3b) at ($(P3)+(0,-0.95)$);
  \coordinate (P4b) at ($(P4)+(0,-0.95)$);
  \coordinate (P5b) at ($(P5)+(0,-0.95)$);
  \coordinate (P6b) at ($(P6)+(0,-0.95)$);
  \coordinate (P7b) at ($(P7)+(0,-0.95)$);
  \coordinate (P8b) at ($(P8)+(0,-0.95)$);

  % 可见的外轮廓、可见棱与前、右底边
  \draw (A) -- (B) -- (P2) -- (P8) -- (P8b) -- (P6b) -- (A);
  \draw (B) -- (P6) -- (P8);
  \draw (P6) -- (P6b);
  \draw (P3) -- (P3b);
  \draw (P7) -- (P7b);

  % 上底面 2×2 网格线（过面心 P4）
  \draw (P3) -- (P5);
  \draw (P1) -- (P7);

  % 被遮挡的棱：后左竖棱、底面后边与左边、三条背面/内部竖棱
  \draw[dashed] (P2) -- (P2b);
  \draw[dashed] (P2b) -- (P8b);
  \draw[dashed] (P2b) -- (A);
  \draw[dashed] (P1) -- (P1b);
  \draw[dashed] (P4) -- (P4b);
  \draw[dashed] (P5) -- (P5b);

  % 顶点标注
  \fill (A) circle (0.85pt) node[below left=-1pt] {$A$};
  \fill (B) circle (0.85pt) node[left=-1pt] {$B$};
  \fill (P1) circle (0.85pt) node[left=-2pt] {$P_1$};
  \fill (P2) circle (0.85pt) node[above left=-1pt] {$P_2$};
  \fill (P3) circle (0.85pt) node[below left=-2pt] {$P_3$};
  \fill (P4) circle (0.85pt) node[above left=-2pt] {$P_4$};
  \fill (P5) circle (0.85pt) node[above=-1pt] {$P_5$};
  \fill (P6) circle (0.85pt) node[below left=-2pt] {$P_6$};
  \fill (P7) circle (0.85pt) node[right=-1pt] {$P_7$};
  \fill (P8) circle (0.85pt) node[above right=-1pt] {$P_8$};
\end{tikzpicture}
